\chapter{Methodology}
\label{chap:methodology}

This chapter formalizes the reconstruction task addressed in the thesis, defines the experimental contrasts that test the hypotheses stated in the introduction, and specifies the architectural, training, and evaluation components used throughout the remaining chapters.


\section{Problem Formulation}
\label{sec:method_problem}

Reservoir simulation produces a sequence of spatial snapshots that describe the thermodynamic state of a fractured-porous medium of Type~II at discrete instants of time. Each snapshot is a three-channel field comprising the global pressure $P$, the water saturation in the fracture network $S$, and the water saturation in the matrix blocks $\bar{S}$. The reconstruction task is formulated at the level of a single time step: for any fixed instant of the simulation, the network receives a coarse-grid realization of these three channels and the auxiliary scalar information available at the same instant, and produces a fine-grid estimate of the same three channels. The temporal evolution is recovered by applying the same reconstruction independently at each time step.

Let $X^{LR} \in \mathbb{R}^{C \times h \times w}$ denote the coarse-grid snapshot with $C = 3$ channels and spatial dimensions $h \times w$, and let $X^{HR} \in \mathbb{R}^{C \times H \times W}$ denote the corresponding fine-grid snapshot with $H = s h$ and $W = s w$, where $s$ is the upscaling factor along each spatial axis. A vector of auxiliary scalars $\mathbf{z} \in \mathbb{R}^{d}$ is associated with the same simulation instant. A super-resolution generator is a parametric mapping
\begin{equation}
    G_{\theta} : \bigl(X^{LR}, \mathbf{z}\bigr) \mapsto \widehat{X}^{HR} \in \mathbb{R}^{C \times H \times W},
    \label{eq:sr_mapping}
\end{equation}
parameterized by $\theta$. When the vector $\mathbf{z}$ is omitted, the mapping reduces to $G_{\theta}(X^{LR})$ and the generator is referred to as unconditional. Training consists in fitting $\theta$ to a finite collection of paired snapshots $\{(X^{LR}_n, X^{HR}_n, \mathbf{z}_n)\}_{n=1}^{N}$ produced by the numerical simulator under varied physical and operational parameters.

The formulation rests on three assumptions. First, the spatial domain is restricted to a two-dimensional vertical cross-section of the reservoir, which is consistent with the geometry of the simulator used to generate the dataset. Second, the upscaling factor $s$ is fixed across all experiments and equal to the ratio between the fine-grid and coarse-grid horizontal resolution; vertical resolution is unchanged between the two grids. Third, the coarse-grid input is an independent numerical solution of the governing equations on a coarse mesh rather than a downsampled version of the fine-grid target, so the mapping that $G_{\theta}$ has to fit recovers physical information lost during coarse-grid simulation rather than inverting a known downsampling operator.

The three hypotheses formulated in the introduction translate into three experimental contrasts that share a common evaluation protocol. The first contrast compares residual generators of different depth, trained with the same pixel-wise objective, on the point-wise reconstruction error of the three-channel field. The second contrast compares the same generators before and after fine-tuning within an adversarial framework, with the comparison performed jointly across point-wise metrics and across structural and spectral metrics, so that the trade-off between the two metric groups is observable. The third contrast compares the deepest residual generator in two configurations, one unconditional and one supplied with the vector $\mathbf{z}$, on the same metric suite.